%% If you are using \orcid or academicons
%% icons, make sure you have the academicons
%% option here, and compile with XeLaTeX
%% or LuaLaTeX.
% \documentclass[10pt,a4paper,academicons]{altacv}

%% Use the "normalphoto" option if you want a normal photo instead of cropped to a circle
% \documentclass[10pt,a4paper,normalphoto]{altacv}

\documentclass[10pt,letter,ragged2e,withhyper]{leozhang}

%% AltaCV uses the fontawesome5 and academicon fonts
%% and packages.
%% See http://texdoc.net/pkg/fontawesome5 and http://texdoc.net/pkg/academicons for full list of symbols. You MUST compile with XeLaTeX or LuaLaTeX if you want to use academicons.

% Change the page layout if you need to
% \geometry{left=1.25cm,right=1.25cm,top=1.5cm,bottom=1.5cm,columnsep=1.2cm}
\geometry{left=1in,right=1in,top=1in,bottom=1in}

% The paracol package lets you typeset columns of text in parallel
\usepackage{paracol}

% Change the font if you want to, depending on whether
% you're using pdflatex or xelatex/lualatex
\ifxetexorluatex
  % If using xelatex or lualatex:
  \setmainfont{Lato}
\else
  % If using pdflatex:
  \usepackage[default]{lato}
\fi

% Change the colours if you want to
\definecolor{VividPurple}{HTML}{3E0097}
\definecolor{SlateGrey}{HTML}{2E2E2E}
\definecolor{LightGrey}{HTML}{666666}
% \colorlet{name}{black}
\colorlet{tagline}{VividPurple}
\colorlet{heading}{VividPurple}
\colorlet{headingrule}{VividPurple}
% \colorlet{subheading}{PastelRed}
\colorlet{accent}{VividPurple}
\colorlet{emphasis}{SlateGrey}
\colorlet{body}{LightGrey}

% Change some fonts, if necessary
% \renewcommand{\namefont}{\Huge\rmfamily\bfseries}
% \renewcommand{\personalinfofont}{\footnotesize}
% \renewcommand{\cvsectionfont}{\LARGE\rmfamily\bfseries}
% \renewcommand{\cvsubsectionfont}{\large\bfseries}

% Change the bullets for itemize and rating marker
% for \cvskill if you want to
\renewcommand{\itemmarker}{{\small\textbullet}}
\renewcommand{\ratingmarker}{\faCircle}

%% sample.bib contains your publications
% \addbibresource{ref.bib}

\newcommand{\etal}{\textit{et al}.}
\newcommand{\ie}{\textit{i}.\textit{e}.}
\newcommand{\eg}{\textit{e}.\textit{g}.}
\renewcommand*{\bibfont}{\small}

\begin{document}

\name{Renyuan Zhang}
\namesuffix{(Leo), Ph.D.}
% \tagline{ADAS Radar Specialist at Lucid Motors}
\tagline{}
% Cropped to square from https://en.wikipedia.org/wiki/Marissa_Mayer#/media/File:Marissa_Mayer_May_2014_(cropped).jpg, CC-BY 2.0
%% You can add multiple photos on the left or right
% \photoR{2.5cm}{mmayer-wikipedia-cc-by-2_0}
% \photoL{2cm}{Yacht_High,Suitcase_High}
\personalinfo{%
  % Not all of these are required!
  % You can add your own with \printinfo{symbol}{detail}
  \email{dr.leo.zhang@outlook.com}
  \phone{669-666-4260}
%   \twitter{@zrmaker}
  \linkedin{zrmaker}
  \github{zrmaker}
  \homepage{Google Scholar}
  %   \mailaddress{Address, Street, 00000 County}
  \location{SF Bay Area}
%   \orcid{orcid.org/0000-0000-0000-0000} % Obviously making this up too. If you want to use this field (and also other academicons symbols), add "academicons" option to \documentclass{altacv}
  %% You MUST add the academicons option to \documentclass, then compile with LuaLaTeX or XeLaTeX, if you want to use \orcid or other academicons commands.
  % \orcid{0000-0000-0000-0000}
  %% You can add your own arbtrary detail with
  %% \printinfo{symbol}{detail}[optional hyperlink prefix]
  % \printinfo{\faPaw}{Hey ho!}
  %% Or you can declare your own field with
  %% \NewInfoFiled{fieldname}{symbol}[optional hyperlink prefix] and use it:
  % \NewInfoField{gitlab}{\faGitlab}[https://gitlab.com/]
  % \gitlab{your_id}
}

\makecvheader

%% Depending on your tastes, you may want to make fonts of itemize environments slightly smaller
\AtBeginEnvironment{itemize}{\small}

\justifying

%% Set the left/right column width ratio to 6:4.
% \columnratio{0.6}
% \setlength{\columnsep}{1.5em}

% Start a 2-column paracol. Both the left and right columns will automatically
% break across pages if things get too long.
% \begin{paracol}{2}

\cvsection{Objective}
Seeking full time engineer/researcher job on ADAS radar system, automotive radar vehicle integration, radar signal processing, radar-related machine learning, sensor fusion and radar target recognition. 


\cvsection{Experience}

\cvevent{ADAS HW Specialist}{\includegraphics[width=50 pt]{icons/lucid.eps} Lucid Motors}{May 2022 -- May 2023}{Newark, CA}
\begin{itemize}
  \item Staff Engineer, component owner of Lucid vehicles Short Range Radar system.
  \item Led the Lucid Gravity Short Range Radar ADAS system architecture defining, signal defining, diagnosis DID/DTC defining, sensor packaging, calibration requirements, ADAS feature defining and requirements, drafted production specifications, assembly work order, done FOV \& KOZ simulations efforts.
  \item Worked with vendor on Lucid Gravity project timeline, SOW, maturity levels, sensor feature deliverables, sensor requirements.
  \item Worked on Lucid Air Short Range Radar sensor system diagnosis, calibration, EU and KSA market retrofit, fascia testing, fleet testing, and service team request.
  \item Led the Lucid Air vendor connections on SW releases, change requests, warranty and FA requests.
  \item Led designing Lucid in-house ADAS radar fascia, bracket, radar testing.
  \item Worked on Lucid level 4 autonomous driving sensor system architecture.
\end{itemize}

\divider

\cvevent{Hardware Engineer}{\includegraphics[width=10 pt]{icons/weride.eps} Weride}{July 2021 -- April 2022}{San Jose, CA}
\begin{itemize}
  \item Product level design of automotive radar. Led the blind spot radar sensor module development using Calterah mmWave AiP chips. 
  \item Developed mmWave radar full signal processing chain using C++ on the programmable MMIC chip.
  \item Involved in in-house radar PCB design and mechanical enclosure radar radome design.
  \item Developed HW-SW data quality control with perception team.
  \item Evaluated new generation radar samples from different vendors and adapted to new POC vehicle sensor builds.
\end{itemize}

\divider

\cvevent{Hardware Engineer}{\faLyft Lyft Inc. (Level 5)}{May 2019 -- May 2021}{Palo Alto, CA}
\begin{itemize}
  \item Led the radar sensor module and submodule characterization, testing and algorithm verification from different vendors.
  \item Proposed algorithms, configured and calibrated radar sensor on for on-vehicle system. Development online calibration for point cloud based automotive radars.
  \item Delivered the radar sensor ERS, IQC and KPIs for verification and functional requirements.
  \item Involved in developing radar early fusion algorithms.
  \item Developed C++ based data quality improvement algorithms for perception including tracking, calibration, interference mitigation and clustering.
  \item Developed radar sensor diagnostic tool for non-Tier 1 grade automotive radars.
  \item Developed binary tools for camera ISP tuning to perception ML outputs.
  \item Delivered part of automated ISP tuning algorithms and automated camera calibration system.
\end{itemize}

\divider

\cvevent{Research Assistant}{\faUniversity University of Arizona}{2015 -- 2019}{Tucson, AZ}
\begin{itemize}
  \item Maintained and developed the TI mmWave ROS interface (\href{https://github.com/radar-lab/ti_mmwave_rospkg}{GitHub link})
  \item Achieved human behavior detection and recognition via CNN using micro-Doppler signatures by mmWave radar.
  \item Developed multi-target multi-input camera-radar sensor fusion and classification.
  \item Developed the mmWave sensing module for hospital ward surveillance including patient behavior recognition, tracking and target posture estimation.
  \item Researched non-synchronized incoherent MIMO radar angle resolution improvements.
  \item Optimized radar point cloud clustering and detection by machine learning method.
  \item Developed CUDA algorithms on radar signal processing.
  \item Researched on radar target clustering and classification.
  \item Realized radar interference detection, classification and mitigation using SVM.
  \item Improved automotive Radar Angular resolution using 3D Luneburg lens.
  \item Completed 3D imaging mmWave circular SAR using single transceiver and compressed sensing.
\end{itemize}



\cvsection{Education}

\cveducation{2015 -- 2019}{Ph.D. in Electrical and Computer Engineering}{University of Arizona}{mmWave Radar: Enhancing Resolution, Target Recognition, and Fusion with Other Sensors}

\divider

\cveducation{2013 -- 2015}{M.S. in Optical Sciences}{University of Arizona}{Fiber Based Spectral Domain Optical Coherence Tomography: Mechanism and Clinical Applications}

\divider

\cveducation{2009 -- 2013}{B.S. in Optoelectronic Engineering}{Chongqing University}{Surface-Enhanced Raman Scattering Substrate Synthesis and Characterization}

% \cvsection{Experience Summary}

% % Adapted from @Jake's answer from http://tex.stackexchange.com/a/82729/226
% % \wheelchart{outer radius}{inner radius}{
% % comma-separated list of value/text width/color/detail}
% % Some ad-hoc tweaking to adjust the labels so that they don't overlap
% \hspace*{-1em}  %% quick hack to move the wheelchart a bit left
% \wheelchart{1.5cm}{0.5cm}{%
%   10/13em/accent!30/Sleeping \& dreaming about work,
%   25/9em/accent!60/Public resolving issues with Yahoo!\ investors,
%   5/11em/accent!10/\footnotesize\\[1ex]New York \& San Francisco Ballet Jawbone board member,
%   20/11em/accent!40/Spending time with family,
%   5/8em/accent!20/\footnotesize Business development for Yahoo!\ after the Verizon acquisition,
%   30/9em/accent/Showing Yahoo!\ \mbox{employees} that their work has meaning,
%   5/8em/accent!20/Baking cupcakes
% }

% Renyuan is an experienced engineer in the field of ADAS and autonomy industries, e.g., Lucid Motors and Lyft autonomous vehicles. He has deep understanding of automotive radar hardware and signal processing, as well as state-of-art algorithms for sensor systems like sensor fusion, radar target recognition, radar waveform design and machine learning for automotive sensors.

% use ONLY \newpage if you want to force a page break for
% ONLY the currentc column
% \newpage


\cvsection{Publications}

\nocite{*}

% \printbibliography[heading=pubtype,title={\printinfo{\faBook}{Books}},type=book]

% \divider


\cvpublications{\faCopy[regular]}{Patents}{}
\begin{itemize}
  \item {\it L269.USNP00}	{\bf R. Zhang}, M. Emadi, A. Mostajeran, J. Izadian, "Sequential Clustering". (Filed 12/27/2019)
  \item {\it L334.USNP00} {\bf R. Zhang}, M. Emadi, A. Mostajeran, J. Izadian, "Blind Online Calibration for Vehicle Radar". (Filed	12/30/2019)
  \item {\it L335.USNP00} {\bf R. Zhang}, M. Emadi, A. Mostajeran, J. Izadian, "Adaptive Clutter Removal From Radar Images". (Filed	12/30/2019)
  \item {\it L338.USNP00} {\bf R. Zhang}, M. Emadi, A. Mostajeran, J. Izadian, "Adaptive Gating for Radar Tracking". (Filed	12/30/2019)
  \item \href{https://patents.justia.com/patent/20190353478}{F. Deng, {\bf R. Zhang}, \etal{}, "System and Method for Angle Measurement,"	United States Patent application 20190353478.}
\end{itemize}
\divider


% \printbibliography[heading=pubtype,title={\printinfo{\faFile*[regular]}{Journal Articles}}, type=article]
\cvpublications{\faFile*[regular]}{Journal Articles}{}
\begin{itemize}
  \item \href{https://ieeexplore.ieee.org/document/8844649}{{\bf R. Zhang} and S. Cao, "Extending Reliability of mmWave Radar Tracking and Detection via Fusion with Camera," {\it IEEE Access}, vol. 7, pp. 137065-137079, 2019.}
  \item \href{https://ieeexplore.ieee.org/document/9083948}{A. Sengupta, {\bf R. Zhang}, \etal{}, "mm-Pose: Real-Time Human Skeletal Posture Estimation using mmWave Radars and CNNs," {\it IEEE Sensors Journal}, vol. 20, no. 17, pp. 10032-10044, 2020.}
  \item \href{https://ieeexplore.ieee.org/document/8585077}{{\bf R. Zhang} and S. Cao, "Real-time Human Motion Behavior Detection via CNN using mmWave Radar," {\it IEEE Sensors Letters}, vol. 3, no. 2, pp. 1-4, 2019.}
  \item \href{http://www.mdpi.com/1424-8220/17/6/1419}{{\bf R. Zhang} and S. Cao, "3D Imaging Millimeter Wave Circular Synthetic Aperture Radar," {\it Sensors}, vol. 17, no. 6, pp. 1419, June 2017.}
\end{itemize}
\divider

% \printbibliography[heading=pubtype,title={\printinfo{\faUsers}{Conference Proceedings}},type=inproceedings]
\cvpublications{\faUsers}{Conference Proceedings}{}
\begin{itemize}
  \item \href{https://ieeexplore.ieee.org/abstract/document/9058256}{{\bf R. Zhang} and S. Cao, "Non-Synchronized Integration using Multiple Radars via Least Squares Fitting," {\it 2019 IEEE National Aerospace and Electronics Conference (NAECON)}, Dayton, OH, 2019, pp. 694-697.}
  \item \href{https://ieeexplore.ieee.org/abstract/document/9048869}{{\bf R. Zhang} and S. Cao, "Robust and Adaptive Radar Elliptical Density-Based Spatial Clustering and Labelling for mmWave Radar Point Cloud Data," {\it 2019 53rd Asilomar Conference on Signals, Systems, and Computers}, Pacific Grove, CA, pp. 919-924.}
  \item \href{https://ieeexplore.ieee.org/document/8835656}{F. Jin, {\bf R. Zhang}, \etal{}, "Multiple Patients Behavior Detection in Real-time using mmWave Radar and Deep CNNs," {\it 2019 IEEE Radar Conference}, Boston, MA, 2019, pp. 1-6.}
  \item \href{https://ieeexplore.ieee.org/abstract/document/8378586}{{\bf R. Zhang} and S. Cao, "Support vector machines for classification of automotive radar interference," {\it 2018 IEEE Radar Conference (RadarConf18)}, Oklahoma City, OK, 2018, pp. 0366-0371.}
  \item \href{http://ieeexplore.ieee.org/document/7944286}{{\bf R. Zhang} and S. Cao, "Compressed Sensing For Portable Millimeter Wave 3D Imaging Radar," {\it 2017 IEEE Radar Conference (RadarConf)}, Seattle, WA, 2017, pp. 0663-0668.}
  \item \href{http://ieeexplore.ieee.org/abstract/document/7944216}{{\bf R. Zhang} and S. Cao, "Portable Millimeter Wave 3D Imaging Radar," {\it 2017 IEEE Radar Conference (RadarConf)}, Seattle, WA, 2017, pp. 0298-0303.}
\end{itemize}

% \cvsection{Life Philosophy}
% \begin{quote}
% ``If you don't have any shadows, you're not standing in the light.''
% \end{quote}

%% Switch to the right column. This will now automatically move to the second
%% page if the content is too long.
% \switchcolumn

\cvsection{Skills}

\cvtag{ADAS}
\cvtag{Automotive Radar}
\cvtag{Sensor System}
\cvtag{Autonomy}

\divider

\cvtag{C/C++}
\cvtag{Python}
\cvtag{Matlab}
\cvtag{Shell}
\cvtag{Source Control}
\cvtag{Algorithm Development}
\cvtag{Data Analysis}

\divider

\cvtag{FMCW}
\cvtag{MIMO}
\cvtag{DSP}
\cvtag{Radar DSP}
\cvtag{Waveform Design}
\cvtag{Image Processing}
\cvtag{Computer Vision}
\cvtag{Sensor Networking}
\cvtag{Sensor Fusion}
\cvtag{Machine Learning}
\cvtag{Denoising}
\cvtag{Target Recognition}

% \newpage


\cvsection{Miscellaneous}

\cvachievement{\faTrophy}{2020 IEEE Sensors Letters Best Paper Award}{Real-Time Human Motion Behavior Detection via CNN Using mmWave Radar}

\divider

\cvachievement{\faGithub}{Public Repo}{The TI mmWave ROS interface which is used by many mmWave radar research groups worldwide (\href{https://github.com/radar-lab/ti_mmwave_rospkg}{GitHub})}

\divider

\cvachievement{\faComments}{Reviewer}{\vspace{-.8\topsep}\begin{itemize}[nosep,topsep=0pt]
    \item IEEE Transactions on Signal Processing
    \item IEEE Sensors Journal 
    \item IEEE Radar Conference
\end{itemize}}

% \divider

% \cvachievement{\faChartLine}{Google's Growth}{from a hundred thousand searches per day to over a billion}

% \newpage

% \cvsection{Referees}

% % \cvref{name}{email}{mailing address}
% \cvref{Prof.\ Alpha Beta}{Institute}{a.beta@university.edu}
% {Address Line 1\\Address line 2}

% \divider

% \cvref{Prof.\ Gamma Delta}{Institute}{g.delta@university.edu}
% {Address Line 1\\Address line 2}

% \end{paracol}

\end{document}
